\documentclass[12pt]{article}

\usepackage{amsmath,amssymb,amsthm}
\usepackage[margin=1in]{geometry}

\theoremstyle{definition}
\newtheorem{definition}{Definition}
\newtheorem{lemma}{Lemma}
\newtheorem{theorem}{Theorem}
\newtheorem{corollary}{Corollary}

\DeclareMathOperator{\lc}{lc}
\newcommand{\Qx}{\mathbb{Q}[x]}

\makeatletter
\renewenvironment{proof}[1][\proofname]{%
  \par\pushQED{\qed}%
  \normalfont \topsep6\p@\@plus6\p@\relax
  \trivlist
  \item[\hskip\labelsep\bfseries #1\@addpunct{.}]\ignorespaces
}{%
  \popQED\endtrivlist\@endpefalse
}
\makeatother

\begin{document}

\begin{lemma}[NIBZO]
There is no integer strictly between $0$ and $1$.  Equivalently, if
$n\in\mathbb{Z}$ and
\[
0<n<1,
\]
then such an $n$ cannot exist.
\end{lemma}

\begin{proof}
Suppose, for contradiction, that there is an integer strictly between $0$ and
$1$.  Put
\[
S=\{n\in\mathbb{Z}:0<n<1\}.
\]
Then $S$ is nonempty and $S\subseteq\mathbb{Z}_{\geq0}$.  By the well-ordering
principle, $S$ has a least element.  Call it $m$.

Since $0<m$, multiplicative closure of positive integers gives $0<m^2$.  Since
$m<1$ and $m>0$, multiplying the inequality $m<1$ by the positive integer $m$
gives
\[
m^2<m.
\]
Thus $m^2$ is also an integer satisfying
\[
0<m^2<m<1.
\]
So $m^2\in S$, but $m^2<m$.  This contradicts the choice of $m$ as the least
element of $S$.  Therefore $S$ is empty.
\end{proof}

\noindent This integer fact matters for $\Qx$ because degrees are integers.  We
are not trying to order the polynomials themselves.  We only compare the degrees
of nonzero polynomials.

\begin{lemma}[No zero divisors in $\mathbb{Q}$]
If $a,b\in\mathbb{Q}$ and $ab=0$, then $a=0$ or $b=0$.
\end{lemma}

\begin{proof}
Suppose $ab=0$.  If $a=0$, there is nothing to prove.  Now assume $a\neq0$.
Then $a^{-1}$ exists in
$\mathbb{Q}$.  Multiplying both sides of $ab=0$ by $a^{-1}$ gives
\[
a^{-1}(ab)=a^{-1}\cdot0.
\]
By associativity, the left side is
\[
(a^{-1}a)b=1b=b.
\]
Also $a^{-1}\cdot0=0$: indeed,
\[
a^{-1}\cdot0=a^{-1}(0+0)=a^{-1}\cdot0+a^{-1}\cdot0,
\]
and adding the additive inverse of $a^{-1}\cdot0$ to both sides gives
$a^{-1}\cdot0=0$.  Hence $b=0$.  Therefore $a=0$ or $b=0$.
\end{proof}

\section{The Ring $\Qx$}

\begin{definition}[The polynomial ring $\Qx$]
We define $\Qx$ to be the set of all finite polynomial expressions
\[
f(x)=a_0+a_1x+\cdots+a_nx^n,
\]
where $n\in\mathbb{Z}_{\geq0}$ and each coefficient $a_i$ lies in $\mathbb{Q}$.
Here finite means that after some power of $x$, all coefficients are zero.
\end{definition}

\begin{definition}[Addition and multiplication]
Let
\[
f(x)=a_0+a_1x+\cdots+a_Nx^N
\qquad\text{and}\qquad
g(x)=b_0+b_1x+\cdots+b_Nx^N
\]
be polynomials in $\Qx$.  If one polynomial originally has smaller degree than
the other, we add zero coefficients so both expressions end at $x^N$.
Their sum is defined coefficient-wise:
\[
f(x)+g(x)=(a_0+b_0)+(a_1+b_1)x+\cdots+(a_N+b_N)x^N.
\]

For multiplication, if
\[
f(x)=a_0+a_1x+\cdots+a_mx^m
\qquad\text{and}\qquad
g(x)=b_0+b_1x+\cdots+b_nx^n,
\]
then
\[
f(x)g(x)=c_0+c_1x+\cdots+c_{m+n}x^{m+n},
\]
where each coefficient $c_k$ is found by adding all products $a_ib_j$ for which
the exponents satisfy $i+j=k$.  For example,
\[
c_0=a_0b_0,\qquad c_1=a_0b_1+a_1b_0,\qquad
c_{m+n}=a_mb_n.
\]
\end{definition}

\begin{definition}[Zero, one, and negatives]
The zero element of $\Qx$ is the zero polynomial
\[
0_{\Qx}=0+0x+0x^2+\cdots.
\]
The one element of $\Qx$ is the constant polynomial
\[
1_{\Qx}=1+0x+0x^2+\cdots.
\]
If
\[
f(x)=a_0+a_1x+\cdots+a_nx^n,
\]
then
\[
-f(x)=(-a_0)+(-a_1)x+\cdots+(-a_n)x^n.
\]
\end{definition}

\begin{theorem}[$\Qx$ is a commutative ring with identity]
With the operations defined above, $\Qx$ is a commutative ring with
identity.  Also,
\[
0_{\Qx}\neq 1_{\Qx}.
\]
\end{theorem}

\begin{proof}
We check the ring conditions.  Addition stays inside $\Qx$ because each
coefficient $a_i+b_i$ is rational.  Multiplication also stays inside $\Qx$
because each coefficient $c_k$ is built from finitely many sums and products of
rational numbers.

The zero polynomial is the additive identity because adding zero to every
coefficient changes nothing.  The polynomial $-f(x)$ is the additive inverse of
$f(x)$ because each coefficient becomes $a_i+(-a_i)=0$.  The constant polynomial
$1$ is the multiplicative identity because multiplying by $1$ leaves the
coefficients unchanged.

The associative and commutative laws come from the same laws in $\mathbb{Q}$,
coefficient by coefficient.  For instance, addition is associative because
\[
(a_k+b_k)+c_k=a_k+(b_k+c_k)
\]
for each coefficient.  The multiplication laws and distributivity are checked in
the same way from the rational-number laws.

Finally, the zero polynomial has constant coefficient $0$, while the constant
polynomial $1$ has constant coefficient $1$.  Since $0\neq1$ in $\mathbb{Q}$,
have $0_{\Qx}\neq1_{\Qx}$.
\end{proof}

\section{Basic Arithmetic Facts}

\begin{lemma}[Uniqueness of zero and one]
The additive identity and multiplicative identity in $\Qx$ are unique.
In particular, if $f(x)+z(x)=f(x)$ for every $f(x)\in\Qx$, then
$z(x)=0$.  Similarly, if $f(x)u(x)=f(x)$ for every $f(x)\in\Qx$, then
$u(x)=1$.
\end{lemma}

\begin{proof}
Suppose $z$ acts like an additive identity for $f$, so $f+z=f$.  Add $-f$ to
both sides:
\[
(-f)+(f+z)=(-f)+f.
\]
By associativity and the definition of additive inverse, this becomes
\[
0+z=0.
\]
Since $0$ is the additive identity, $z=0$.

Now suppose $u$ acts like a multiplicative identity for every polynomial.  Test
this on the constant polynomial $1$:
\[
1\cdot u=1.
\]
But $1\cdot u=u$.  Hence $u=1$.
\end{proof}

\begin{lemma}[Uniqueness of additive inverses]
If $u(x)$ and $v(x)$ are both additive inverses of $f(x)$, then $u(x)=v(x)$.
\end{lemma}

\begin{proof}
Assume
\[
f+u=0
\qquad\text{and}\qquad
f+v=0.
\]
Then $f+u=f+v$.  Add $-f$ to both sides:
\[
(-f)+(f+u)=(-f)+(f+v).
\]
By associativity,
\[
((-f)+f)+u=((-f)+f)+v.
\]
Thus
\[
0+u=0+v,
\]
so $u=v$.
\end{proof}

\begin{lemma}[Zero and negatives]
For all $f(x),g(x)\in\Qx$,
\[
f(x)\cdot0=0,\qquad -(-f(x))=f(x),\qquad (-f(x))(-g(x))=f(x)g(x).
\]
\end{lemma}

\begin{proof}
First,
\[
f\cdot0=f(0+0)=f\cdot0+f\cdot0.
\]
Adding the additive inverse of $f\cdot0$ to both sides gives $f\cdot0=0$.

Next, since
\[
(-f)+f=0,
\]
the polynomial $f$ is an additive inverse of $-f$.  By uniqueness of additive
inverses,
\[
-(-f)=f.
\]

For the product of negatives, first observe that
\[
fg+(-f)g=(f+(-f))g=0\cdot g=0.
\]
Therefore $(-f)g=-(fg)$.  Applying this with $g$ replaced by $-g$ gives
\[
(-f)(-g)=-(f(-g)).
\]
Also,
\[
fg+f(-g)=f(g+(-g))=f\cdot0=0,
\]
so $f(-g)=-(fg)$.  Hence
\[
(-f)(-g)=-(-(fg))=fg.
\]
\end{proof}

\section{Degree and the Integral Domain Property}

\begin{definition}[Degree and leading coefficient]
If
\[
f(x)=a_0+a_1x+\cdots+a_nx^n
\]
with $a_n\neq0$, then the degree of $f$ is
\[
\deg(f)=n.
\]
The leading coefficient is
\[
\lc(f)=a_n.
\]
The zero polynomial has no degree.  When convenient, we write
$\deg(0)=-\infty$ only as notation.
\end{definition}

\begin{lemma}[The integer measure for $\Qx$]
If $f(x)\in\Qx$ is nonzero, then $\deg(f)$ is an integer in
$\mathbb{Z}_{\geq0}$.  Therefore integer facts such as NIBZO and the
well-ordering principle apply to polynomial degrees.
\end{lemma}

\begin{proof}
By definition, a nonzero polynomial has the form
\[
f(x)=a_0+a_1x+\cdots+a_nx^n
\]
with $a_n\neq0$ and $n\in\mathbb{Z}_{\geq0}$.  Its degree is defined to be this
integer $n$.  Hence $\deg(f)\in\mathbb{Z}_{\geq0}$.
\end{proof}

\begin{lemma}[Degree of a product]
If $f(x),g(x)\in\Qx$ are nonzero, then
\[
\deg(fg)=\deg(f)+\deg(g).
\]
\end{lemma}

\begin{proof}
Let the leading term of $f$ be $a_mx^m$ and the leading term of $g$ be
$b_nx^n$, where $a_m,b_n\neq0$.  The highest possible power in $fg$ is
$x^{m+n}$.  The only way to get this power is to multiply $a_mx^m$ by
$b_nx^n$, so the coefficient of $x^{m+n}$ is $a_mb_n$.  Since $\mathbb{Q}$ has
no zero divisors by the lemma above, $a_mb_n\neq0$.  Therefore $fg$ has degree
$m+n$.
\end{proof}

\begin{theorem}[$\Qx$ has no zero divisors]
If $f(x),g(x)\in\Qx$ and
\[
f(x)g(x)=0,
\]
then $f(x)=0$ or $g(x)=0$.
\end{theorem}

\begin{proof}
Suppose $f\neq0$ and $g\neq0$.  Then
\[
\deg(fg)=\deg(f)+\deg(g)\geq0,
\]
so $fg$ is nonzero.  This contradicts $fg=0$.  Hence at least one of $f$ and
$g$ must be zero.
\end{proof}

\begin{corollary}[Cancellation]
If $h(x)\neq0$ and
\[
h(x)f(x)=h(x)g(x),
\]
then $f(x)=g(x)$.
\end{corollary}

\begin{proof}
Subtract the right side from the left:
\[
h(f-g)=0.
\]
Since $h\neq0$ and $\Qx$ has no zero divisors, we must have
$f-g=0$.  Thus $f=g$.
\end{proof}

\section{Divisibility}

\begin{definition}[Divisibility]
For $a(x),b(x)\in\Qx$, we say $a(x)$ divides $b(x)$, written
\[
a(x)\mid b(x),
\]
if there exists $q(x)\in\Qx$ such that
\[
b(x)=a(x)q(x).
\]
\end{definition}

\section{The Division Algorithm in $\Qx$}

For $\Qx$, degree plays the role that size plays in integer division.  Since
degrees are nonnegative integers, we can use the well-ordering principle to pick
a remainder of smallest degree.

\begin{theorem}[Division algorithm in $\Qx$]
Let $f(x),g(x)\in\Qx$ with $g(x)\neq0$.  Then there exist unique
polynomials $q(x),r(x)\in\Qx$ such that
\[
f(x)=q(x)g(x)+r(x),
\]
where either
\[
r(x)=0
\]
or
\[
\deg(r)<\deg(g).
\]
\end{theorem}

\begin{proof}
\textbf{Existence.}
Look at all possible remainders after subtracting multiples of $g$ from $f$:
\[
S=\{\,f(x)-q(x)g(x):q(x)\in\Qx\,\}.
\]
If $0\in S$, then $f=qg$ for some $q$, so
\[
f=qg+0,
\]
and the result holds with $r=0$.

Now suppose $0\notin S$.  Then every element of $S$ is nonzero, so
each element has a degree in $\mathbb{Z}_{\geq0}$.  Define
\[
D=\{\,\deg(s):s\in S\,\}\subseteq\mathbb{Z}_{\geq0}.
\]
The set $D$ is nonempty because $f-0g=f$ lies in $S$.  By the well-ordering
principle, $D$ has a smallest element.  Choose $r\in S$ with minimum degree.
Since $r\in S$, there exists $q\in\Qx$ such that
\[
r=f-qg.
\]
Equivalently,
\[
f=qg+r.
\]

It remains to show that $\deg(r)<\deg(g)$.  Suppose instead that
\[
\deg(r)\geq\deg(g).
\]
Write
\[
r(x)=a_mx^m+\cdots
\qquad\text{and}\qquad
g(x)=b_nx^n+\cdots,
\]
where $a_m,b_n\neq0$ and $m\geq n$.  Since $b_n\neq0$ and $b_n\in\mathbb{Q}$,
the rational number $a_m/b_n$ exists.  Define
\[
r^*(x)=r(x)-\frac{a_m}{b_n}x^{m-n}g(x).
\]
The second term has leading term $a_mx^m$, so the leading term of $r$ cancels.
Thus either $r^*=0$ or
\[
\deg(r^*)<\deg(r).
\]
Also,
\[
\begin{aligned}
r^*
&=f-qg-\frac{a_m}{b_n}x^{m-n}g\\
&=f-\left(q+\frac{a_m}{b_n}x^{m-n}\right)g.
\end{aligned}
\]
Therefore $r^*\in S$.  If $r^*=0$, then $0\in S$, which is not the case we are
in.  If $r^*\neq0$, then $r^*$ is in $S$ and has smaller degree than $r$,
contradicting how we chose $r$.  Therefore our assumption was false, and
\[
\deg(r)<\deg(g).
\]

\textbf{Uniqueness.}
Suppose
\[
f=q_1g+r_1=q_2g+r_2,
\]
where $r_1$ and $r_2$ are both either zero or have degree less than $\deg(g)$.
Subtracting gives
\[
g(q_1-q_2)=r_2-r_1.
\]
Assume $q_1\neq q_2$.  Then $q_1-q_2\neq0$, and since $\Qx$ has no
zero divisors,
\[
g(q_1-q_2)\neq0.
\]
Hence
\[
\deg(g(q_1-q_2))=\deg(g)+\deg(q_1-q_2)\geq\deg(g).
\]
But $r_2-r_1$ is either zero or has degree strictly less than $\deg(g)$, since
both $r_1$ and $r_2$ have degree less than $\deg(g)$.  This is impossible.
Therefore $q_1=q_2$.  Substituting back gives $r_1=r_2$.
\end{proof}

\end{document}
